\section{Discussion}
\label{sec:discussion}

\textbf{The central result of this study is not that confidence-guided
allocation has been shown to outperform an equal uniform allocation of
the same total resources.} The matched-resource controls do not support
that claim. What the evidence shows is that the proposed rules can infer
and deploy a workable reconstruction budget from the observed measurement
data, whereas the original pipeline requires that budget to be selected
manually in advance. Every statement elsewhere in this article is scoped
to that reading.

Three findings sit uneasily together at first reading: the combination
of all three modules removes essentially every failure in the regime
where failures occur; individually, no module beats a budget-matched
control; and those controls are themselves not implementable, since they
require knowing in advance what the adaptive rule would have chosen. The
reconciliation is that the modules are not competing with a well-chosen
fixed budget --- they are competing with the problem of choosing one.
Secs.~\ref{sec:ofat} and~\ref{sec:combined} showed that no single starved
budget causes failure while three starved together cause it half the time, so a
practitioner cannot diagnose which knob to turn from a failure alone, and
the correct setting varies cell by cell. That the rule's allocation
performs like an oracle-matched uniform one is the strongest available
evidence that it picks sensible values.

Table~\ref{tab:adopt} states what we would recommend a practitioner
adopt, and on what evidence. The modules are \emph{not} equally
supported, and presenting them as a uniform package would misrepresent
the study.

\begin{table}[htbp]
\centering\small
\caption{Adoption guidance. ``Evidence'' summarises
Secs.~\ref{sec:results}--\ref{sec:matched}; costs are from the companion
article~\cite{methodpaper}.}
\label{tab:adopt}
\begin{tabular}{@{}lp{0.36\linewidth}p{0.36\linewidth}@{}}
\toprule
Module & Evidence & Recommendation \\
\midrule
\textbf{B} coverage
 & Largest and most stable effect; improved every instance with
   headroom; broke at most two trials anywhere.
 & \textbf{Adopt first.} Cheapest of the three: no additional circuit
   execution, same asymptotic class as the code it replaces. \\[3pt]
\textbf{D} shot stopping
 & Large where shots bind, nil where they do not; the most two-sided
   per-trial behaviour in the study.
 & Adopt where sampling is cheap relative to a failed run; it dominates
   wall-clock cost whenever it fires. \\[3pt]
\textbf{C} beam width
 & Fails Holm correction; null on two of three held-out instances;
   indistinguishable from its matched control on every trial.
 & \textbf{Do not enable alone.} Retained in the Full configuration only
   because it provably never narrows the beam and 20 attributed failures
   needed both budget axes relaxed together. \\
\bottomrule
\end{tabular}
\end{table}

None of this is in tension with the calibration result reported in the
companion article~\cite{methodpaper}: module D's stopping rule needs
$\Cw$ to order windows correctly, not to state a correct probability. Where the
miscalibration does cost something is the choice of $\varepsilon$, which
consequently has no principled setting and was left at its default.

% =====================================================================

% =====================================================================
